\section{평가}
\label{sec:eval}

본 장에서는 무방어 구현(\textsf{baseline})과 세 가지 대응 기법---전체 이중 연산(\textsf{double}),
선행 연구~\cite{lee2026haetae}의 대응을 재구현한 \textsf{leeha}, 그리고 제안 기법
\textsf{IRV}---을 동일한 조건에서 (i)~정확성, (ii)~커버리지, (iii)~2차 오류 내성,
(iv)~구현 비용의 네 가지 측면에서 비교한다.

\subsection{정확성}
\begin{itemize}
  \item 무결함 시 세 대응 기법 모두 \textsf{baseline}과 \textbf{동일한 서명(동일 다이제스트)}을 방출한다.
  \item \textsf{IRV}는 잔차 $\delta=0$이 보장되어 거짓 양성(오탐)이 발생하지 않는다.
\end{itemize}

\subsection{커버리지}
\label{subsec:coverage}
각 (오류, 변형) 쌍에 대해 오류를 주입하고, 방출된 서명의 다이제스트를 golden(무영향),
\textsc{leak}(\textsf{baseline}과 동일하여 미차단), blocked(무효화)의 세 가지로 분류한다
(표~\ref{tab:coverage}).

\begin{table}[H]\centering
\caption{오류 커버리지 및 2차 오류 내성. \texttt{blk}=blocked, \textsc{leak}=미차단, 우회=탐지 분기
스킵으로 무력화. $\dagger$ = 본 연구 도입 지점.}
\label{tab:coverage}
\small
\begin{tabular}{@{}lcccc@{}}
\toprule
오류 지점 & baseline & \textsf{double} & \textsf{leeha} & \textsf{IRV} \\
\midrule
SEED / SIGNBIT / UNPACK / LSB & \textsc{leak} & blk & blk & blk \\
T1 (CS)$^\dagger$ / T2 (ADDY)$^\dagger$ & \textsc{leak} & blk & blk & blk \\
RB (REJECT)$^\dagger$ & \textsc{leak} & blk & \textbf{\textsc{leak}} & blk \\
\midrule
차단 합계 / 7 & 0 & \textbf{7} & 6 & \textbf{7} \\
\midrule
2차 오류 (T2 + 탐지 분기 스킵) & \textsc{leak} & \textbf{\textsc{leak}(우회)} & \textbf{\textsc{leak}(우회)} & \textbf{차단} \\
\bottomrule
\end{tabular}
\end{table}

\textsf{double}과 \textsf{IRV}는 7개 지점을 모두 차단한 반면, \textsf{leeha}는 6개 지점을
차단하되 \textbf{거부 우회(RB)만 차단하지 못하였다}. RB는 본 논문의 위협 모델에 속하는 지점으로
선행 연구~\cite{lee2026haetae}의 범위 밖이다. 따라서 이 결과는 해당 연구의 오류로 해석해서는 안
되며, ``서명 후 검증 기법이 $B_1<B_2$로 인해 RB를 구조적으로 놓친다''는 의미로 읽어야 한다
(\ref{subsec:b1b2}절). 한편 \textsf{baseline}의 T2 \textsc{leak}이 실제 비밀키 누설임은
\ref{subsec:key-recovery}절에서 100\% 키 복원으로 입증한 바 있다.

\subsection{2차 오류 내성}
\label{subsec:2nd}
1차 오류(T2, 누설 생성)와 2차 오류(대응 기법의 탐지 분기 스킵)를 함께 주입하여, 탐지 분기를
무력화하는 공격에 대한 내성을 평가한다(표~\ref{tab:coverage} 마지막 행). 이는 표~\ref{tab:coverage}의
RB 커버리지(1차 오류 차단)와는 별개인 \emph{2차} 오류 위협에 대한 평가이다.

\textsf{double}의 결과 비교 분기와 \textsf{leeha}의 검증 분기는 모두 단일 스킵으로 우회되어, 누설
서명이 \textsf{baseline}과 동일하게 그대로 방출되었다. 반면 IRV는 응답 연산의 핵심 지점(T1/T2/RB)을
비교 분기 없이 감염형으로 처리하므로, 건너뛸 분기 자체가 존재하지 않아 차단을 유지하였다(상류 지점은
표준 검증으로 추가 보호된다). 다만 본 평가의 2차 오류는 \emph{탐지 분기 스킵} 모델에 한하며, 감염 연산
자체(마스크 적용$\cdot$잔차 누적)를 표적으로 하는 2차 오류는 \ref{sec:disc}절에서 별도로 논한다.

\subsection{비용}
\label{subsec:cost}
무결함 빌드의 코드/RAM(\texttt{arm-none-eabi-size})과 서명 1회 사이클(DWT)을 측정한다
(표~\ref{tab:cost}).

\begin{table}[H]\centering
\caption{구현 비용(무결함 빌드). 시간은 baseline 대비 상대값.}
\label{tab:cost}
\small
\begin{tabular}{@{}lccc@{}}
\toprule
변형 & 서명시간 & 코드(text) & 정적 RAM \\
\midrule
baseline       & $1.00\times$ & 24{,}556\,B & 31{,}168\,B \\
\textsf{double}& $2.00\times$ & $+0.7\%$    & $+0.6\%$ \\
\textsf{leeha} & $1.14\times$ & $+25.4\%$   & $+20.0\%$ \\
\textsf{IRV}   & $1.14\times$ & $+26.3\%$   & $+20.7\%$ \\
\bottomrule
\end{tabular}
\end{table}

이중 연산은 서명 시간이 2배로 증가하는 반면, \textsf{IRV}와 \textsf{leeha}는 각각 약 $1.14\times$에
그쳤다(표준 검증 1회와 상수 크기 검사를 포함). 코드 크기 증가 약 $25\%$는 주로 연결된 검증
루틴($\sim$6\,KB)에서 비롯된다. 즉 \textbf{서명 시간} 기준으로 IRV는 이중 연산의 $60\%$ 미만 비용으로
전면 커버리지와 거부 우회 차단, 2차 오류 내성을 동시에 달성한다. 다만 \textbf{코드/정적 RAM} 측면에서는
이중 연산(코드 $+0.7\%$, RAM $+0.6\%$)이 IRV($+26.3\%$, $+20.7\%$)보다 작으므로, IRV와 이중 연산은
(서명 시간$\cdot$2차 오류 내성) 대 (코드$\cdot$RAM)의 \emph{트레이드오프} 관계로 이해해야 한다.

\paragraph{\textsf{leeha} 비용에 관한 주석.}
공정한 비교를 위해 본 재구현의 \textsf{leeha}는 채택 이후 부분 이중 검사와 검증을 방출 서명당 한
번씩 수행하였으며, 이에 따라 $+14\%$의 비용이 발생하였다. 원 논문~\cite{lee2026haetae}은 더 가벼운
부분 이중 연산($+0.4\%$)과 서명 후 검증($+4.6\%$)을 보고하므로, 본 논문은 두 수치를 함께 제시한다.
다만 비용에 관한 핵심 결론은 이 세부 사항과 무관하다. 즉 IRV는 이중 연산의 절반 미만 비용으로
전면 커버리지와 거부 우회 차단, 2차 오류 내성을 모두 확보한다.

\subsection{종합}
\label{subsec:summary}
\begin{table}[H]\centering
\small
\begin{tabular}{@{}lcccl@{}}
\toprule
 & 커버리지 & RB & 2차오류 & 비용(시간 ; 코드/RAM) \\
\midrule
\textsf{double} & 7/7 & $\checkmark$ & 우회 & $2.0\times$ ; $+0.7\%/+0.6\%$ \\
\textsf{leeha}  & 6/7 & $\times$ & 우회 & $1.14\times$ ; $+25.4\%/+20.0\%$ \\
\textsf{IRV}    & \textbf{7/7} & \textbf{$\checkmark$} & \textbf{생존} & $1.14\times$ ; $+26.3\%/+20.7\%$ \\
\bottomrule
\end{tabular}
\end{table}

IRV는 (1)~거부 우회를 포함한 7개 지점 전면 차단, (2)~2차 오류 생존, (3)~이중 연산 절반 수준의
\emph{서명 시간}을 동시에 만족하는 유일한 변형이다. 두 대응 기법과의 비교를 나누어 보면 다음과 같다.
\begin{itemize}
  \item \textbf{vs \textsf{double}}: 커버리지는 7/7로 동일하나, IRV는 (i)~서명 시간이 절반($1.14\times$ 대
        $2.0\times$ --- IRV는 검증 1회, 이중 연산은 서명 전체를 2번 계산)이고, (ii)~2차 오류에 생존한다
        (이중 연산의 결과 비교 분기는 단일 스킵으로 우회된다). 대신 코드$\cdot$정적 RAM은 \textsf{double}이
        더 작아, 이 부분은 트레이드오프이다.
  \item \textbf{vs \textsf{leeha}}: 서명 시간은 비슷($\approx1.14\times$)하나, IRV는 (i)~거부 우회(RB)를
        차단하고 (ii)~2차 오류에 생존한다는 점에서 앞선다 --- \textsf{leeha}는 둘 다 놓친다(원 논문의
        경량 구성은 \ref{subsec:cost}절 주석 참조).
\end{itemize}

\subsection{논의 및 한계}
\label{sec:disc}

본 절에서는 제안 기법과 실험의 한계를 명확히 하고, 이를 둘러싼 해석상의 주의점을 정리한다.

\begin{itemize}
  \item \textbf{선행 기법의 차용.} IRV의 표준 검증(M6), 부분 이중 검사(M3), 정상성 검사(M2)는
        선행 연구~\cite{lee2026haetae}와 탐지 원리를 공유한다. 본 논문의 기여는 새로운 탐지 항목을
        제시한 데 있지 않고, 이들을 \emph{비교 연산 없는 감염형으로 통합}하여 2차 오류 내성을 부여하고,
        서명 경계 재검사로 거부 우회를 차단하며, 단일 서명 T2 공격을 오류 모델 하에서 실증한 데 있다.
  \item \textbf{거부 우회의 범위.} RB의 \textsc{leak}은 ``서명 경계를 위반한 응답이 검증을
        통과한다''는 의미이며, T2와 같은 즉각적 키 복원이 아니라 통계적 누설에 해당한다.
  \item \textbf{\textsf{leeha} 비용.} \ref{subsec:cost}절의 수치는 공정 비교를 위한 재구현 결과이며,
        원 논문의 보고치와 함께 제시하였다.
  \item \textbf{오류 모델의 현실성.} 본 논문이 가정한 명령어 스킵은 32비트 마이크로컨트롤러에서
        물리적으로 실증된 표준 오류 모델이며~\cite{moro2013}, 소프트웨어 라인 주입은 이 오류의
        \emph{논리적} 효과를 결정론적으로 재현하여 대응 기법을 공정하게 비교한다. 본 논문의 공격 및
        대응 평가는 이 에뮬레이션된 단일 오류 모델을 기준으로 하며, 물리적 오류 주입을 통한 실현은
        향후 과제로 남긴다.
  \item \textbf{2차 오류 위협 모델.} 본 논문은 출력을 게이팅하는 \emph{조건 분기} 1회 스킵을
        모델링하며, IRV의 생존은 응답 연산을 비교 분기 없이 처리한 데서 비롯된다. 이 모델은
        detect-and-abort의 구조적 약점(뒤집을 불리언 분기의 존재)을 정확히 겨냥한 것으로, IRV는
        그러한 분기를 제거한다. 다만 동일 능력의 공격자는 감염 적용 단계(마스크 XOR)나 잔차 누적
        자체를 표적으로 하는 2차 오류를 시도할 수 있으며, 이 경우 감염형 대응도 무력화될 수 있다.
        즉 본 논문의 2차 오류 내성 주장은 \emph{탐지 분기 스킵 모델에 한정}되며, 감염 연산을
        데이터 흐름에 직조하여 이러한 잔존 위협까지 견디게 하는 강건 구현은 향후 과제다.
  \item \textbf{대리 주입과 결정론 설정.} 소프트웨어 라인 주입과 결정론 난수 발생기는 대응 기법을
        공정하게 비교하기 위한 설정이며, 복구 절차는 소프트웨어 클린 스킵 누설에 대한 self-test로 100\% 검증하였다.
  \item \textbf{절대 수치의 키 의존성.} \textsf{baseline}의 60.86\,M 사이클은 해당 키의 거부 반복
        횟수에 의존하므로, 본 논문은 \emph{상대} 비용을 기준으로 결론을 제시한다.
  \item \textbf{일반화.} 제안 원리는 HAETAE 120/180/260은 물론, 동일한 응답 구조를 갖는 Fiat--Shamir
        with aborts 계열(\texttt{Dilithium}~\cite{dilithium} 등)로 확장할 수 있다.
\end{itemize}
